\documentclass[11pt]{article}
\usepackage[utf8]{inputenc}
\usepackage{geometry}
\geometry{a4paper, margin=1in}
\usepackage{titlesec}
\usepackage{enumitem}
\usepackage{xcolor}

\title{\textbf{Strategic Diagnostic Snapshot™}}
\author{RankPilot Intelligence}
\date{\today}

\begin{document}

\maketitle

\section*{I. Executive Summary}
\textbf{Firm:} González de Araujo Consultores \\
\textbf{Practice Area:} Data Protection \\
\textbf{Detected Practice Model:} Standard

\section*{II. Structural Positioning}
\textbf{Current Tier Assessment:} Unranked / New Entry \\
\textbf{Strategic Confidence:} 100.0\% \\
\textbf{Advantage:} Standard Market Offerings

\section*{III. Portfolio Analysis (Key Matters)}
\begin{itemize}
\item \textbf{Grupo Hermes Data Protection Framework} (Client: Grupo Hermes | Value: N/A)\\ Grupo Hermes engaged the firm to design and implement a thorough data protection and governance framework, placing the team in a central role in translating personal-data obligations into defined internal controls and operational procedures. The engagement was driven by the need to establish an organised approach to the handling of personal data, including the identification of how that information moved through the client’s operations and the creation of documentation and processes tailored to those data flows.

Arturo González de Araujo and Ernesto Zavala led the work. The team mapped Grupo Hermes’s personal data flows, providing the factual foundation for the framework’s implementation. Building on that mapping exercise, the firm developed tailored privacy notices, implemented internal policies governing the handling of personal data, and established procedures for ARCO rights management. These were concrete deliverables directed at the full lifecycle of ARCO requests: creating an internal process through which requests could be managed consistently and in accordance with the framework designed for Grupo Hermes.

The stated outcome was 100% compliance in handling ARCO requests and avoiding regulatory sanctions. This result links the firm’s work on data-flow mapping, tailored privacy notices, internal policies and ARCO rights-management procedures to a measurable compliance outcome for Grupo Hermes. For the submission narrative, the matter evidences Arturo González de Araujo’s and Ernesto Zavala’s role in designing and implementing data protection governance that operates through specific, client-facing processes rather than policy documentation alone.
\item \textbf{Mega Direct Data Protection Advisory} (Client: MEGA DIRECT, S.A. de C.V. | Value: N/A)\\ Mega Direct Data Protection Advisory concerns the firm’s 16-year relationship with MEGA DIRECT, S.A. de C.V. in connection with complex data protection and database-related matters. The duration of this engagement is material: rather than reflecting a discrete instruction, it records continuing work over 16 years on issues affecting the client’s handling of data and the operation of its databases.

The team’s role extended beyond general data protection advice. The firm represented MEGA DIRECT, S.A. de C.V. in legal proceedings concerning data acquisition, processing, storage, and protection. This required the team to address each of those distinct stages of the data-handling cycle—acquisition, processing, storage, and protection—in the context of proceedings affecting the client’s business activities.

Arturo González de Araujo, Ernesto Zavala, and Damián Ortega were the lead partners on the engagement. Their work combined the long-term advisory dimension of the 16-year relationship with representation in the relevant legal proceedings. The stated client impact was the preservation of business continuity in a highly regulated environment, where the treatment and protection of data and the management of database-related issues were directly connected to the client’s ongoing operations.

For the submission, the matter provides evidence of sustained data protection and database-related work over 16 years, coupled with proceedings representation concerning data acquisition, processing, storage, and protection. It demonstrates a long-duration institutional engagement in which Arturo González de Araujo, Ernesto Zavala, and Damián Ortega supported MEGA DIRECT, S.A. de C.V. on matters bearing directly on business continuity.
\item \textbf{Biocodex Data Protection Framework} (Client: Biocodex | Value: N/A)\\ Biocodex Data Protection Framework

Biocodex engaged the team to design and implement a structured personal data protection framework. The engagement was directed at moving from a general need for personal-data governance to a defined internal structure capable of supporting the client’s management of personal-data responsibilities. The work was not limited to a conceptual assessment: it encompassed both the design of the framework and its implementation within Biocodex’s internal organisation.

Arturo González de Araujo, Ernesto Zavala and Damián Ortega led the engagement for Biocodex. Their work supported the establishment of a dedicated Personal Data Protection Department, giving the personal data protection function a distinct internal home and a clear basis for the application of the structured framework. By combining framework design with implementation work and the creation of the dedicated department, the team addressed the client’s internal organisation of personal-data protection as an operational function.

The stated result was the strengthening of Biocodex’s internal controls and a reduction in regulatory exposure. This matter provides concrete evidence of the team’s ability to translate personal data protection requirements into institutional arrangements within a client organisation: designing and implementing a structured personal data protection framework while helping establish a dedicated Personal Data Protection Department to support stronger internal controls and reduced regulatory exposure.
\item \textbf{Hotel Riazor Data Protection Enhancement} (Client: Hotel Riazor | Value: N/A)\\ Hotel Riazor engaged the team on a focused data-protection mandate directed at strengthening its personal data protection and governance framework. The engagement centred on the practical and institutional question of how Hotel Riazor handles personal data lawfully: not simply as a matter of policy drafting, but through the internal structures responsible for applying personal-data requirements on an ongoing basis. The work therefore addressed both the client’s governance arrangements and the operational support required by its Personal Data Department.

Arturo González de Araujo, Ernesto Zavala and Damián Ortega led the firm’s work for Hotel Riazor. Their role was to support the client in reinforcing its Personal Data Department and in implementing internal policies for the lawful handling of personal data. By pairing departmental reinforcement with the implementation of internal policies, the team helped Hotel Riazor place responsibility for personal-data handling within a more clearly supported internal governance framework.

The result of the engagement was the reinforcement of Hotel Riazor’s Personal Data Department and the implementation of internal policies intended to ensure lawful handling of personal data. For the submission narrative, the matter provides specific evidence of the team’s work on the operationalisation of personal data protection and governance: strengthening the internal department charged with these responsibilities while implementing policies designed to govern the client’s handling of personal data lawfully.
\item \textbf{Grupo Excelsior Data Protection and Compliance} (Client: Grupo Excelsior | Value: N/A)\\ Grupo Excelsior instructed the team on the regularisation and restructuring of its operations, with data protection as the central focus of the engagement. The mandate required the team to translate data-protection requirements into practical operational measures within the client’s regularisation and restructuring exercise. This was not limited to high-level commentary: the work addressed how the client’s internal operations would be organised around data-protection considerations as part of the wider restructuring process.

Arturo González de Araujo, Ernesto Zavala and Damián Ortega led the matter for Grupo Excelsior. Their role was to design and implement the internal policies and training programmes required for the client’s restructured operations. The internal policies established documented standards for the treatment of data-protection issues within the business, while the training programmes equipped personnel to apply those standards in the course of their operational responsibilities. The combination of policy design, implementation and training connected the legal requirements of data protection directly to the client’s regularisation process.

The result was a significant strengthening of Grupo Excelsior’s compliance posture and a reduction in its regulatory exposure. The matter demonstrates the team’s ability to undertake data-protection work through an operational regularisation and restructuring mandate, delivering internal policies and training programmes that support the practical implementation of compliance requirements. Arturo González de Araujo, Ernesto Zavala and Damián Ortega’s involvement evidences a three-partner team deployment focused on converting data-protection requirements into implemented internal measures for Grupo Excelsior.
\item \textbf{Grupo Modelquipo Data Protection Framework} (Client: Grupo Modelquipo | Value: N/A)\\ Grupo Modelquipo engaged the firm to implement a data protection framework as a component of its corporate governance strategy. The mandate positioned data protection within the client’s corporate governance agenda, connecting the implementation exercise directly to the organisation’s approach to risk management and compliance rather than treating personal-data issues as a standalone exercise.

Arturo González de Araujo, Ernesto Zavala and Damián Ortega led the engagement. The team’s work focused on implementing the data protection framework in a manner aligned with Grupo Modelquipo’s corporate governance strategy. In particular, the mandate brought together the framework’s practical implementation with the risk-management and compliance considerations identified by the client, ensuring that data protection was addressed through the structures and culture relevant to corporate governance.

The stated result was an improvement in regulatory compliance and the integration of data protection into Grupo Modelquipo’s governance culture. This is significant because the engagement evidences a governance-led approach to data protection: Arturo González de Araujo, Ernesto Zavala and Damián Ortega supported Grupo Modelquipo in connecting implementation, risk management and compliance with the client’s wider corporate governance strategy. The matter therefore demonstrates the team’s ability to translate data protection requirements into an integrated governance-culture initiative with a defined compliance objective.
\item \textbf{Tiendas Chedraui Data Protection Advisory} (Client: Tiendas Chedraui | Value: N/A)\\ Tiendas Chedraui engages the firm in connection with the legal assessment required for new store openings and its operational expansion. The mandate places data protection at the point at which new consumer-facing operations are being introduced and scaled, requiring the client to consider the legal consequences of the proposed operating model alongside its expansion decisions. In a high-volume consumer data environment, the opening of new stores and the extension of operational activity create a material need to identify the data-processing activities associated with those developments and to evaluate the risks arising from them.

Arturo González de Araujo and Damián Ortega lead the engagement. Their role is to carry out the legal assessment of the data-processing risks linked to Tiendas Chedraui’s new store openings and operational expansion, and to assess whether the relevant processing arrangements comply with applicable regulatory requirements. The work therefore addresses the practical interface between the client’s operational growth and the regulatory obligations governing consumer data, focusing on the risk profile created when high volumes of consumer data are processed in connection with expanded operations.

The significance of the matter lies in its direct connection to Tiendas Chedraui’s expansion activity: the firm evaluates risks associated with data processing in a high-volume consumer data environment while ensuring compliance with regulatory requirements. The mandate evidences an ability to translate data-protection requirements into a legal assessment tailored to concrete new store openings and operational expansion, rather than addressing consumer-data compliance only at a general level.

\end{itemize}

\section*{IV. Strategic Evolution Path}
\begin{itemize}
\item Maintain current strategy.
\end{itemize}

\vfill
\centering
\small{This document is a structural assessment and does not guarantee ranking results.}

\end{document}